\documentclass[a4paper,12pt]{ctexbook}

\usepackage{amsmath}
\usepackage{amssymb}
\usepackage{amsthm}
\usepackage{geometry}
\geometry{left=2cm, right=2cm, top=3cm, bottom=3cm}
\usepackage{setspace}
\onehalfspacing
\usepackage{hyperref}
\hypersetup{hidelinks}

\theoremstyle{definition}
\newtheorem{definition}{定义}[chapter]
\newtheorem{example}{例}[chapter]

\theoremstyle{plain}
\newtheorem{theorem}{定理}[chapter]
\newtheorem{lemma}{引理}[chapter]
\newtheorem{proposition}{命题}[chapter]
\newtheorem{corollary}{推论}[chapter]

\theoremstyle{remark}
\newtheorem*{remark}{注}

\title{格物数言}
\author{一个问题}
\date{}

\begin{document}

\maketitle

\frontmatter

\chapter*{前言}

本书旨在通过一些其他（不来自于课本）的思路来理解（初中）数学，以便让读者无需生啃那些无聊的概念，而是让读者意识到数学是一门美妙的艺术，从而激起学习数学（进而包括其它门类）的兴趣。因此本书的叙述以
通俗性为主，而不是具有很强的逻辑性。

同时本书也不是一本正宗的教材，仅可作为数学爱好者的娱乐读物和激起学习兴趣的敲门砖。

最后感谢我的朋友对于我的支持，特别感谢哀的旋律，她为本书取了一个相当好的名字。

\tableofcontents

\mainmatter

\chapter{函数、方程、不等式的关系}

方程、函数、不等式，这几个看似风马牛不相及的数学工具，其实却有着相同的源头，让我们从本质上来探寻这几个数学工具的关系。

\section{函数}
函数是数学中一个重要的工具，它在不同阶段有不同的定义，接下来给出三种不同阶段的定义：

\begin{definition}
在一个变化过程中，有2个变量$x$和$y$，如果对于每一个自变量$x$，$y$都有唯一的值与之对应，则称$y$是$x$的函数。
\end{definition}

\begin{definition}
设$A$、$B$是非空数集，如果按照某种确定的对应关系，使得对于集合$A$中的任意一个数$x$，在集合$B$中都有唯一确定的数$y$与之对应，那么称$f:A \to B$为集合$A$到集合$B$的一个函数。
\end{definition}

\begin{definition}
设 $A$、$B$ 为非空集合，若二元关系 $f \subseteq A \times B$ 满足：
\begin{enumerate}
\item 存在性：$\forall x \in A,\ \exists y \in B,\ (x,y) \in f$；
\item 唯一性：$\forall x \in A,\ \forall y_1, y_2 \in B,\ 
\bigl((x,y_1) \in f \land (x,y_2) \in f\bigr) \Rightarrow y_1 = y_2$；
\end{enumerate}
则称 $f$ 为从 $A$ 到 $B$ 的一个函数，记作 $f: A \to B$。
\end{definition}

不难看出，它们都描述一种对应（映射）的关系。其中定义1.1是初中的定义，更侧重于描述变化过程中2个变量之间的关系；定义1.2是高中阶段的定义，限定在“数集”的范围内；
定义1.3是最严格的定义，严谨地表述了函数的本质。

讲到这里，就不得不简单提一下函数的来源：为了描述变化。

在没有函数这个概念之前，对于一组数据，我们只能将它们看做是一个个静态的点。但是，随着科学的进步，人们发现单一的静态的结果已经不再满足发展的需要，需要有一种严谨的数学工具来描述变化。
于是函数诞生了。

有了函数这个概念之后，人们用变化的眼光来看这个世界，人们不仅能看到“当下”，更能通过函数来看到“将来”。随着数学的一次又一次进步，函数逐渐变得更加多样。如今，函数已经成为描述生活的不可或缺的工
具。这又与另一个重要的数学工具：微积分，扯上了关系。它们的关系简单来说是这样的：函数是微积分的研究对象，而微分和积分是作用于函数的两种基本预算。

最后在节末附上2句有关函数的名言：

“数学的本质在于它的自由。”——康托尔

“函数是自然的语言，微积分是它的语法。”——费曼（多认为）

\begin{remark}
以下的章节若无说明，则默认在初中定义下讨论。
\end{remark}

\section{不等式}

不等式在不同的阶段有着不同的定义，以下给出两种定义：

\begin{definition}
用不等号连接，表示不等关系的式子叫做不等式。
\end{definition}

\begin{definition}
不等式是一个包含一个或多个未知量的陈述句，该陈述句使用不等号来连接两个数学表达式，并断言这两个表达式之间具有某种大小或顺序关系。这个陈述句的真假取决于未知量在指定范围内所取的具体数值，因此它是一
个真假待定的开语句。
\end{definition}

从以上的两个定义中似乎不大能够看出不等式和函数的关系，那就让我们从变化的角度来理解不等式。

前面讲到函数的起源——描述变化。如果说函数是为了描述变化而诞生的数学工具，那么不等式就是为了描述一段范围而生的数学工具。

让我们把它与函数联系起来，不等式即是描述函数的结果在某一范围内（函数值的取值范围）时，应当具备什么样的条件（自变量的范围）。
反过来也成立：给出一定的条件（自变量的取值范围），就可以判断函数的结果在什么范围内（函数值的取值范围）。

因此我们完全可以说：不等式是函数的的部分（满足一定的条件），因此通过不等式，我们就可以知道函数的范围。

从图象上来讲则更为直观：不等式是函数图象满足一定条件的一段。打个比方，原来的函数是一块肉，在刀（条件）的切割下满足要求的那部分。

\section{方程}

方程和上面的函数，不等式一样，也是有着不同的定义，以下给出2种定义：

\begin{definition}
含有未知数的等式叫做方程。使方程左右两边相等的未知数的值，叫做方程的解。求方程的解的过程，叫做解方程。
\end{definition}

\begin{definition}
方程是一个包含一个或多个未知量的陈述句，该陈述句使用等号连接两个数学表达式并断言两者相等。该陈述句的真假取决于未知量在指定范围内所取的具体数值，因此它是一个真假待定的开语句。在事先给定的未知量取值范围内，将某个具体数值代入方程中的未知量后，能够使该方程所断言的相等陈述句变成一个真命题，则该数值称为该方程的一个解。所有能够满足这一条件的数值组成的整体，称为该方程的解集。依据逻辑推理和等价变形规则，从给定方程出发，在指定论域中寻找并确定其全部解的推演过程。该过程必须保证既不引入非原方程解的值，也不遗漏任何满足原方程的值，最终以解集形式明确输出所有解。这样的过程叫做解方程。
\end{definition}

单看这两种定义，方程似乎和前面的函数与不等式似乎毫无关联，那就让我们再次从函数的角度出发，在初中的范围内，来理解方程究竟是想描述什么。

由以上的研究，我们知道函数实际上就是变量之间的对应关系，而方程就是描述在这个过程中的特殊位置的数学工具。即对于函数的一个确定的值，我们如何求出自变量的值。也就是说，方程就是为了求出在变化过程中的一个特殊位置所需的特殊条件而诞生的数学工具。

因此我们说：方程是函数上的一个“点”。它是求特定条件的工具。

从更高等的数学上讲，方程并不一定是点，但无论如何，它总是比对应的函数低一纬。

从图象上来讲，对于一个简单的函数，其与一条平行于$x$轴的直线（特定条件）的交点就可以列出一个方程，而点的坐标就是这个方程的解。

总而言之，如果将函数看做一棵大树，那么不等式就是一段枝干，方程就是一片树叶。即不等式是函数的部分，方程是函数的一个点。

这样，我们就弄清楚了函数、方程、不等式之间的关系。


\chapter{命题}

命题和证明是数学的基础。这一章节我们将讨论课本只是简略提到的命题与证明。让你简单领略数学推理的魅力。

\section{定义}

在学习数学的过程中，对于一个新概念，我们总是习惯的说，“这个东西的定义是什么？”。有了定义，我们就能清楚的知道数学概念的特点，本质，以及从定义出发，得到许多有用的结论。因此，让我们对“定义”下一个定义（事实上，从数学的角度上来讲，我们并不能给定义做出一个良好的定义）：

\begin{definition}
对于一个概念做出明确规定的语句，叫做这个概念的定义
\end{definition}

一个正确，良好的定义是数学推理的开始。这就是定义的作用：严谨地描述一个新概念。

\section{命题}

\end{document}